\section{Desafíos y Oportunidades de Mejora}

Durante el proceso de implementación y puesta en marcha del sistema Odoo 19 en Machu Picchu
Exclusive Tours E.I.R.L., se identificaron retos operativos críticos que fueron gestionados
para consolidar la transformación digital de la agencia, así como áreas estratégicas para el
crecimiento futuro.

\subsection{Desafíos Identificados durante la Implementación}
\begin{itemize}
    \item \textbf{Gestión de Seguridad de Datos Sensibles:} uno de los mayores retos fue migrar la recepción de fotografías de pasaportes desde dispositivos personales y WhatsApp hacia el repositorio cifrado de Odoo. Anteriormente, por temor a vulnerabilidades de seguridad, la gerencia borraba periódicamente estas imágenes, lo que impedía contar con un historial fiable para las estrictas regulaciones de ingreso a sitios como Machu Picchu.
    \item \textbf{Centralización de Procesos en una Estructura Unipersonal:} debido a que la empresa posee una estructura vertical donde la gerencia asume múltiples roles (desde ventas hasta coordinación logística), el desafío consistió en diseñar una interfaz simplificada que redujera la carga administrativa en un 20\% sin complicar la operación diaria.
    \item \textbf{Estandarización de Itinerarios Variables:} adaptar el sistema para gestionar viajes que oscilan entre 2 y 20 días requirió una configuración detallada en el módulo de Proyectos. El reto fue asegurar que cada etapa del tablero Kanban reflejara la realidad de servicios personalizados que incluyen desde traslados simples hasta paquetes completos con hoteles y vuelos.
    \item \textbf{Control de Riesgo Financiero en Reservas:} se identificó la necesidad crítica de bloquear la compra de servicios no reembolsables (boletos de tren o avión) hasta que el registro manual del adelanto (30\% al 50\%) fuera verificado en el sistema, eliminando la incertidumbre de las "ventas a ciegas".
\end{itemize}

\subsection{Oportunidades de Mejora}
A partir de la estabilidad alcanzada con los módulos de CRM y Proyectos, se han detectado las
siguientes áreas para potenciar la competitividad internacional de la agencia:

\begin{itemize}
    \item \textbf{Integración de Marketing Digital y Captación Directa:} existe una oportunidad inmediata para vincular el CRM con campañas segmentadas en Meta Ads y Google Ads. Esto permitirá reducir la dependencia actual de las recomendaciones orgánicas por WhatsApp y captar activamente leads en los mercados norteamericano y europeo.
    \item \textbf{Automatización de la Fidelización (Post-Venta):} el sistema permite ahora programar actividades automáticas para el seguimiento de clientes tras su viaje, como el envío de tarjetas personalizadas por cumpleaños o fin de año, fortaleciendo el vínculo emocional con el turista.
    \item \textbf{Habilitación de Inteligencia de Negocios (BI):} se recomienda aprovechar los datos acumulados para generar reportes sobre nacionalidades con mayor demanda y rentabilidad por tour. Esta información permitirá a la gerencia optimizar el presupuesto publicitario y ajustar las estrategias de precios según la temporada.
    \item \textbf{Escalabilidad hacia Integraciones Externas:} a futuro, la agencia podría conectar Odoo con plataformas de reservas globales (OTAs) como Booking o Expedia, sincronizando automáticamente las ventas externas con el flujo operativo interno de la empresa.
    \item \textbf{Adopción de Pagos Integrados:} aunque actualmente se usan links de pago externos, la futura integración de una pasarela de pagos directa con Odoo automatizaría el registro de abonos, eliminando totalmente la digitación manual de saldos.
\end{itemize}
